%% TOAN12_C5_ONTAP2.tex -- Đề ôn tập chương V, đề số 2
%% Theo cấu trúc đề thi tốt nghiệp THPT: 12 TN, 4 Đ/S, 6 trả lời ngắn.
%% Đáp án: \dapan{khoá}{hướng dẫn}.
\begin{deontap}{2}
\phanI
\begin{cau}Trong không gian $Oxyz$, điểm nào dưới đây thuộc mặt phẳng $(P)\colon -x+y+z+3=0$?
\pa{$Q(-2;3;-9)$}{$Q(-2;2;-8)$}{$Q(-3;4;-9)$}{$Q(-2;3;-8)$}
\dapan{D}{Thay toạ độ vào phương trình $(P)$.}\end{cau}
\begin{cau}Trong không gian $Oxyz$, cho mặt cầu $(S)\colon x^2+(y-2)^2+(z-2)^2=16$ và điểm $M(-3;3;1)$. Khẳng định nào sau đây đúng? Điểm $M$
\pa{là tâm của $(S)$}{thuộc $(S)$}{nằm ngoài $(S)$}{nằm trong $(S)$}
\dapan{D}{$IM^2=11$ so với $R^2=16$.}\end{cau}
\begin{cau}Trong không gian $Oxyz$, góc giữa hai đường thẳng $\Delta\colon \begin{cases}x=1+t\\y=0\\z=-1\end{cases}$ và $\Delta'\colon \begin{cases}x=s\\y=2+s\\z=1\end{cases}$ bằng
\pa{$45^\circ$}{$60^\circ$}{$30^\circ$}{$90^\circ$}
\dapan{A}{$\cos(\Delta,\Delta')=\dfrac{|1|}{1\cdot\sqrt{2}}$ nên góc bằng $45^\circ$.}\end{cau}
\begin{cau}Trong không gian $Oxyz$, đường thẳng $\Delta\colon \dfrac{x-1}{-1}=\dfrac{y}{-1}=\dfrac{z+3}{2}$ có một vectơ chỉ phương là
\pa{$\vec u=(-1;1;2)$}{$\vec u=(-1;-1;2)$}{$\vec u=(-1;1;3)$}{$\vec u=(1;0;-3)$}
\dapan{B}{Đọc mẫu số của phương trình chính tắc.}\end{cau}
\begin{cau}Trong không gian $Oxyz$, mặt cầu tâm $I(2;1;2)$ và tiếp xúc với mặt phẳng $(P)\colon 2x+2y-z+3=0$ có phương trình là
\pa{$(x-2)^2+(y-1)^2+(z-2)^2=\dfrac{49}{9}$}{$(x-2)^2+(y-1)^2+(z-2)^2=\dfrac{7}{3}$}{$(x-2)^2+(y-1)^2+(z-2)^2=7$}{$(x+2)^2+(y+1)^2+(z+2)^2=\dfrac{49}{9}$}
\dapan{A}{$R=\mathrm d(I,(P))=\dfrac{7}{3}$.}\end{cau}
\begin{cau}Trong không gian $Oxyz$, góc giữa đường thẳng $\Delta$ có vectơ chỉ phương $\vec u=(0;0;1)$ và mặt phẳng $(P)\colon x+z-3=0$ bằng
\pa{$0^\circ$}{$30^\circ$}{$60^\circ$}{$45^\circ$}
\dapan{D}{$\sin(\Delta,(P))=\dfrac{1}{\sqrt{2}}$ nên góc bằng $45^\circ$.}\end{cau}
\begin{cau}Trong không gian $Oxyz$, cho hai điểm $A(-1;-3;-2)$ và $B(-2;-2;-2)$. Mặt phẳng trung trực của đoạn thẳng $AB$ có phương trình là
\pa{$x-y-1=0$}{$x-y=0$}{$x-y-2=0$}{$3x+5y+4z+25=0$}
\dapan{A}{Đi qua trung điểm $I(- \dfrac{3}{2};- \dfrac{5}{2};-2)$, vectơ pháp tuyến $\vt{AB}=(-1;1;0)$: $x-y-1=0$.}\end{cau}
\begin{cau}Trong không gian $Oxyz$, cho hai đường thẳng $\Delta_1\colon \begin{cases}x=1+t\\y=2-t\\z=-1+2t\end{cases}$ và $\Delta_2\colon \begin{cases}x=3+s\\y=2+s\\z=1\end{cases}$. Vị trí tương đối của $\Delta_1$ và $\Delta_2$ là
\pa{Cắt nhau}{Trùng nhau}{Chéo nhau}{Song song}
\dapan{A}{$\vec u_1=(1;-1;2)$, $\vec u_2=(1;1;0)$, $[\vec u_1,\vec u_2]=(-2;2;2)$, $\vt{A_1A_2}\cdot[\vec u_1,\vec u_2]=0$.}\end{cau}
\begin{cau}Một trạm phát sóng đặt tại $T(3;5;2)$ trong hệ trục $Oxyz$ (đơn vị: km), vùng phủ sóng là khối cầu bán kính $4$ km. Một người dùng điện thoại tại $M(-2;7;6)$. Khoảng cách (km) từ người đó đến trạm (làm tròn đến hàng phần mười) là
\hinh{\begin{tikzpicture}[scale=0.75,font=\footnotesize,line join=round,line cap=round,>=Stealth]
\draw[thick] (-3.2,0)--(3.2,0); \fill[pattern=north east lines] (-3.2,-.15) rectangle (3.2,0);
\draw[line width=1.5pt] (0,0)--(0,.6); \draw (-.2,0)--(0,.6)--(.2,0); \fill (0,.6) circle (.05) node[right]{$T$};
\draw[dashed] (0,.6) circle (2.4); \draw[dashed] (-2.33,0) arc (180:360:2.33 and .3); \draw[<->] (0,.6)--(1.7,2.3) node[midway,above left]{$R$};
\filldraw[fill=gray!40] (2.6,0) rectangle (2.9,.35); \node[above] at (2.75,.35){nhà};
\end{tikzpicture}}

\pa{$45$}{$2{,}7$}{$13{,}4$}{$6{,}7$}
\dapan{D}{$TM=\sqrt{45}\approx 6{,}7$ (km) $>$ $4$ nên người đó ở ngoài vùng phủ sóng.}\end{cau}
\begin{cau}Trong không gian $Oxyz$ (đơn vị: km), một máy bay cất cánh từ điểm $A(-37;37;0)$ và bay theo đường thẳng có vectơ chỉ phương $\vec u=(4;2;1)$. Khi máy bay đạt độ cao $20$ km (tức $z=20$) thì nó ở vị trí $M(a;b;20)$. Giá trị $a+b$ bằng
\hinh{\begin{tikzpicture}[scale=0.85,font=\footnotesize,line join=round,line cap=round,>=Stealth]
\draw[->,gray] (0.000,0.000)--(-1.440,-1.080) node[below left]{$x$};
\draw[->,gray] (0.000,0.000)--(5.200,0.000) node[right]{$y$};
\draw[->,gray] (0.000,0.000)--(0.000,3.400) node[above]{$z$};
\node[below right,gray] at (0.000,0.000){$O$};
\fill[gray!10] (-0.300,0.000)--(-1.596,-0.972)--(3.504,-0.972)--(4.800,0.000)--cycle;
\draw[thick] (-0.336,0.298)--(3.768,2.476); \draw[->,thick] (-0.336,0.298)--(1.716,1.387);
\fill (-0.336,0.298) circle (.05) node[below right]{$A$}; \fill (3.768,2.476) circle (.05) node[above]{$B$};
\draw[dashed] (-0.336,0.298)--(-0.336,-0.702) (3.768,2.476)--(3.768,-0.324);

\end{tikzpicture}}

\pa{$120$}{$20$}{$-34$}{$86$}
\dapan{A}{Đường bay: $\begin{cases}x=-37+4t\\ y=37+2t\\ z=t\end{cases}$. $z=20\Leftrightarrow t=20$, $M(43;77;20)$, $a+b=120$.}\end{cau}
\begin{cau}Trong không gian $Oxyz$ (đơn vị: km), một máy bay cất cánh từ điểm $A(43;42;0)$ và bay theo đường thẳng có vectơ chỉ phương $\vec u=(3;6;1)$. Khi máy bay đạt độ cao $10$ km (tức $z=10$) thì nó ở vị trí $M(a;b;10)$. Giá trị $a+b$ bằng
\hinh{\begin{tikzpicture}[scale=0.85,font=\footnotesize,line join=round,line cap=round,>=Stealth]
\draw[->,gray] (0.000,0.000)--(-1.440,-1.080) node[below left]{$x$};
\draw[->,gray] (0.000,0.000)--(5.200,0.000) node[right]{$y$};
\draw[->,gray] (0.000,0.000)--(0.000,3.400) node[above]{$z$};
\node[below right,gray] at (0.000,0.000){$O$};
\fill[gray!10] (-0.300,0.000)--(-1.596,-0.972)--(3.504,-0.972)--(4.800,0.000)--cycle;
\draw[thick] (-0.336,0.298)--(3.768,2.476); \draw[->,thick] (-0.336,0.298)--(1.716,1.387);
\fill (-0.336,0.298) circle (.05) node[below right]{$A$}; \fill (3.768,2.476) circle (.05) node[above]{$B$};
\draw[dashed] (-0.336,0.298)--(-0.336,-0.702) (3.768,2.476)--(3.768,-0.324);

\end{tikzpicture}}

\pa{$175$}{$-29$}{$185$}{$146$}
\dapan{A}{Đường bay: $\begin{cases}x=43+3t\\ y=42+6t\\ z=t\end{cases}$. $z=10\Leftrightarrow t=10$, $M(73;102;10)$, $a+b=175$.}\end{cau}
\begin{cau}Một rađa đặt tại gốc toạ độ có vùng quét là khối cầu bán kính $20$ km. Một máy bay bay ở độ cao $10$ km theo mặt phẳng $z=10$. Vùng máy bay có thể bị phát hiện là hình tròn có bán kính (km) bằng
\hinh{\begin{tikzpicture}[scale=0.85,font=\footnotesize,line join=round,line cap=round,>=Stealth]
\draw[thick] (-2.8,0)--(2.8,0); \fill[pattern=north east lines] (-2.8,-.15) rectangle (2.8,0);
\draw[dashed] (2.4,0) arc (0:180:2.4); \fill (0,0) circle (.06) node[below]{rađa}; \draw[thick] (-2.7,1.4)--(2.7,1.4) node[right]{đường bay};
\filldraw[fill=gray!50] (1,1.4)--(1.5,1.4)--(1.6,1.5)--(.95,1.5)--cycle; \draw[<->] (-2.3,0)--(-2.3,1.4) node[midway,left]{$h$};
\draw[thick] (-1.91,1.4)--(1.91,1.4); \draw (0,0)--(1.91,1.4) node[midway,below right]{$R$};
\end{tikzpicture}}

\pa{$20$}{$20 \sqrt{3}$}{$10 \sqrt{3}$}{$10$}
\dapan{C}{Mặt phẳng $z=10$ cách tâm $10$ km, cắt mặt cầu theo đường tròn bán kính $r=\sqrt{20^2-10^2}=10 \sqrt{3}$ (km).}\end{cau}
\phanII
\begin{cau}Trong không gian $Oxyz$, cho ba điểm $A(2;0;0)$, $B(0;2;0)$, $C(0;0;1)$.
\begin{dungsai}
  \y Mặt phẳng $(ABC)$ có phương trình $2x+2y+z=1$. \oDS
  \y Vectơ $\vec n=(2;2;1)$ là một vectơ pháp tuyến của $(ABC)$. \oDS
  \y Khoảng cách từ $O$ đến $(ABC)$ bằng $\dfrac{\sqrt{6}}{3}$. \oDS
  \y Thể tích khối tứ diện $OABC$ bằng $\dfrac{2}{3}$. \oDS
\end{dungsai}
\dapan{S,S,Đ,Đ}{$(ABC)\colon x+y+2z-2=0$; $\mathrm d(O,(ABC))=\dfrac{\sqrt{6}}{3}$; $V=\dfrac{2\cdot2\cdot1}6$.}\end{cau}
\begin{cau}Trong không gian $Oxyz$, cho đường thẳng $\Delta\colon \begin{cases}x=-1-t\\y=1-t\\z=-t\end{cases}$.
\begin{dungsai}
  \y Điểm $(-2;0;-1)$ thuộc $\Delta$. \oDS
  \y $\Delta$ song song với mặt phẳng $(P)\colon -x-y-z-3=0$. \oDS
  \y $\Delta$ cắt $(P)$ tại điểm $(-2;0;-1)$. \oDS
  \y Vectơ $(-1;-1;1)$ là một vectơ chỉ phương của $\Delta$. \oDS
\end{dungsai}
\dapan{Đ,S,Đ,S}{$t=1$ cho $(-2;0;-1)$; $\vec u$ cùng phương $\vec n_P$; thay vào $(P)$: $t=1$.}\end{cau}
\begin{cau}Trong không gian $Oxyz$, cho mặt cầu $(S)\colon x^2+y^2+z^2-4z-5=0$.
\begin{dungsai}
  \y Mặt cầu $(S)$ có tâm $I(0;0;-2)$. \oDS
  \y Mặt cầu $(S)$ có bán kính $R=3$. \oDS
  \y Điểm $M(3;0;2)$ thuộc $(S)$. \oDS
  \y Mặt phẳng $(Oxy)$ cắt mặt cầu $(S)$. \oDS
\end{dungsai}
\dapan{S,Đ,Đ,Đ}{$I(0;0;2)$, $R=3$; $\mathrm d(I,(Oxy))=2$.}\end{cau}
\begin{cau}Một đường dốc cho xe lăn được mô hình hoá trong hệ trục $Oxyz$ (đơn vị: m): chân dốc $A(0;0;0)$, đỉnh dốc $B(15;0;2)$, mặt đất là $(Oxy)$. Xét tính đúng sai của các mệnh đề sau:
\hinh{\begin{tikzpicture}[scale=0.85,font=\footnotesize,line join=round,line cap=round,>=Stealth]
\fill[gray!25] (0,0)--(5,0)--(5,.9)--cycle; \draw[thick] (0,0)--(5,.9)--(5,0)--cycle;
\node[below left] at (0,0){$A$}; \node[above right] at (5,.9){$B$}; \draw (.9,0) arc (0:10.2:.9); \node at (1.25,.12){$\alpha$};
\draw[<->] (5.25,0)--(5.25,.9) node[midway,right]{độ cao}; \draw[<->] (0,-.3)--(5,-.3) node[midway,below]{chiều dài theo phương ngang};
\end{tikzpicture}}

\begin{dungsai}
  \y Góc giữa mặt dốc và mặt đất xấp xỉ $82{,}4^\circ$. \oDS
  \y Sin của góc giữa mặt dốc và mặt đất bằng $\dfrac{2 \sqrt{229}}{229}$. \oDS
  \y Độ dài mặt dốc $AB$ bằng $\sqrt{229}$ m. \oDS
  \y Theo quy chuẩn, đường dốc cho xe lăn có độ dốc (tỉ số chiều cao trên chiều dài nằm ngang) không vượt quá $\dfrac{1}{12}$. Đường dốc này đạt quy chuẩn. \oDS
\end{dungsai}
\dapan{S,Đ,Đ,S}{$\overrightarrow{AB}=(15;0;2)$, mặt đất có vectơ pháp tuyến $(0;0;1)$, $\sin\alpha=\dfrac{2}{\sqrt{229}}$. a) Sai: đó là phần bù của góc cần tìm. b) Đúng. c) Đúng. d) Sai: vì $\dfrac{2}{15}>\dfrac{1}{12}$.}\end{cau}
\phanIII
\begin{cau}Trong không gian $Oxyz$, đường thẳng $\Delta\colon \begin{cases}x=-2t\\y=-2-2t\\z=2t\end{cases}$ cắt mặt phẳng $(P)\colon x-y+z+1=0$ tại điểm $M(a;b;c)$. Tính $a+b+c$.\otraloi
\dapan{1}{$t=- \dfrac{3}{2}$, $M(3;1;-3)$.}\end{cau}
\begin{cau}Trong không gian $Oxyz$, mặt phẳng $(P)$ đi qua ba điểm $A(2;0;0)$, $B(0;3;0)$, $C(0;0;6)$. Tính thể tích khối tứ diện $OABC$.\otraloi
\dapan{6}{$V=\dfrac{2\cdot3\cdot6}{6}=6$.}\end{cau}
\begin{cau}Trong không gian $Oxyz$, gọi $M'(a;b;c)$ là điểm đối xứng của điểm $M(-1;2;0)$ qua mặt phẳng $(P)\colon x+2y+2z+6=0$. Tính $a+b+c$.\otraloi
\dapan{-9}{Hình chiếu $H(-2;0;-2)$ là trung điểm $MM'$; $M'(-3;-2;-4)$.}\end{cau}
\begin{cau}Trong không gian $Oxyz$, gọi $M'(a;b;c)$ là điểm đối xứng của điểm $M(-1;0;2)$ qua mặt phẳng $(P)\colon x+2y+2z-12=0$. Tính $a+b+c$.\otraloi
\dapan{11}{Hình chiếu $H(0;2;4)$ là trung điểm $MM'$; $M'(1;4;6)$.}\end{cau}
\begin{cau}Trong một hệ trục $Oxyz$ (đơn vị: mét), một tấm kính phẳng được đặt đi qua ba điểm $A(8;0;0)$, $B(0;5;0)$, $C(0;0;4)$. Một đèn nhỏ đặt tại điểm $E(6;4;1)$. Khoảng cách (mét) từ đèn đến tấm kính bằng bao nhiêu (làm tròn kết quả đến hàng phần trăm)?
\hinh{\begin{tikzpicture}[scale=0.85,font=\footnotesize,line join=round,line cap=round,>=Stealth]
\draw[->,gray] (0.000,0.000)--(-1.260,-0.945) node[below left]{$x$};
\draw[->,gray] (0.000,0.000)--(4.300,0.000) node[right]{$y$};
\draw[->,gray] (0.000,0.000)--(0.000,3.200) node[above]{$z$};
\node[below right,gray] at (0.000,0.000){$O$};
\fill[gray!25,opacity=.8] (-0.936,-0.702)--(3.400,0.000)--(0.000,2.500)--cycle; \draw[thick] (-0.936,-0.702) node[left]{$A$}--(3.400,0.000) node[below]{$B$}--(0.000,2.500) node[left]{$C$}--cycle;
\shade[ball color=gray!20] (0.640,0.030) circle (.1); \node[right] at (0.640,0.030){$E$};
\end{tikzpicture}}
\otraloi
\dapan{2{,}33}{Mặt phẳng $(ABC)\colon \dfrac{x}{8}+\dfrac{y}{5}+\dfrac{z}{4}=1\Leftrightarrow 5x+8y+10z-40=0$. $d=\dfrac{|128|}{\sqrt{3024}}\approx 2{,}33$.}\end{cau}
\begin{cau}Trong không gian $Oxyz$ (đơn vị: km), một máy bay cất cánh từ điểm $A(-26;-32;0)$ và bay theo đường thẳng có vectơ chỉ phương $\vec u=(3;-4;1)$. Khi máy bay đạt độ cao $20$ km (tức $z=20$) thì nó ở vị trí $M(a;b;20)$. Giá trị $a+b$  bằng bao nhiêu?
\hinh{\begin{tikzpicture}[scale=0.85,font=\footnotesize,line join=round,line cap=round,>=Stealth]
\draw[->,gray] (0.000,0.000)--(-1.440,-1.080) node[below left]{$x$};
\draw[->,gray] (0.000,0.000)--(5.200,0.000) node[right]{$y$};
\draw[->,gray] (0.000,0.000)--(0.000,3.400) node[above]{$z$};
\node[below right,gray] at (0.000,0.000){$O$};
\fill[gray!10] (-0.300,0.000)--(-1.596,-0.972)--(3.504,-0.972)--(4.800,0.000)--cycle;
\draw[thick] (-0.336,0.298)--(3.768,2.476); \draw[->,thick] (-0.336,0.298)--(1.716,1.387);
\fill (-0.336,0.298) circle (.05) node[below right]{$A$}; \fill (3.768,2.476) circle (.05) node[above]{$B$};
\draw[dashed] (-0.336,0.298)--(-0.336,-0.702) (3.768,2.476)--(3.768,-0.324);

\end{tikzpicture}}
\otraloi
\dapan{-78}{Đường bay: $\begin{cases}x=-26+3t\\ y=-32-4t\\ z=t\end{cases}$. $z=20\Leftrightarrow t=20$, $M(34;-112;20)$, $a+b=-78$.}\end{cau}
\end{deontap}
